\documentclass[11pt]{article}
\usepackage[utf8]{inputenc}
\usepackage{geometry}
\geometry{a4paper, margin=1in}
\usepackage{graphicx}
\usepackage{amsmath}
\usepackage{amssymb}
\usepackage{amsfonts}
\usepackage{booktabs}
\usepackage{float}
\usepackage{hyperref}
\usepackage{listings}
\usepackage{xcolor}
\usepackage{caption}
\usepackage{subcaption}
\usepackage{enumitem}
\usepackage{appendix}

% Define colors for code listings
\definecolor{codegreen}{rgb}{0,0.6,0}
\definecolor{codegray}{rgb}{0.5,0.5,0.5}
\definecolor{codepurple}{rgb}{0.58,0,0.82}
\definecolor{backcolour}{rgb}{0.95,0.95,0.92}

\lstdefinestyle{mystyle}{
    backgroundcolor=\color{backcolour},
    commentstyle=\color{codegreen},
    keywordstyle=\color{magenta},
    numberstyle=\tiny\color{codegray},
    stringstyle=\color{codepurple},
    basicstyle=\ttfamily\footnotesize,
    breakatwhitespace=false,
    breaklines=true,
    captionpos=b,
    keepspaces=true,
    numbers=left,
    numbersep=5pt,
    showspaces=false,
    showstringspaces=false,
    showtabs=false,
    tabsize=2
}
\lstset{style=mystyle}

\title{\textbf{Project 1: Co-creativity as Curation} \\ \large MovieLens 100K Ranking \& Personalization Pipeline}
\author{
    Team \#[Insert Team Number] \\
    [Member 1 Name] \\
    [Member 2 Name] \\
    [Member 3 Name] \\
    [Member 4 Name] \\
    [Member 5 Name]
}
\date{\today}

\begin{document}

\maketitle

\begin{abstract}
This report presents the implementation and evaluation of an end-to-end co-creative curation pipeline for movie recommendation on the MovieLens 100K dataset. The system encompasses candidate generation, ranking, reranking with diversity constraints, and a personalization update mechanism. We implement and compare multiple ranking approaches: a popularity baseline, Matrix Factorization trained with Stochastic Gradient Descent (MF-SGD), Matrix Factorization trained with Alternating Least Squares (MF-ALS), and a pairwise Learning-to-Rank (LTR) model. Furthermore, we incorporate Maximal Marginal Relevance (MMR) reranking to explore the relevance-diversity tradeoff. A personalization update using Exponential Moving Average (EMA) is demonstrated to simulate session-based user state evolution. Offline evaluation using Recall@10, NDCG@10, and Diversity@10 reveals the comparative strengths of each method and the impact of MMR. The results indicate that while MF-based methods improve relevance, the LTR model achieves the best ranking performance. MMR successfully increases diversity at a controllable cost to relevance.
\end{abstract}

\section{Introduction}
\label{sec:introduction}

The goal of this project is to build a co-creative curation system that generates a personalized top-K slate of movies for a user. Following the system blueprint outlined in the assignment, our pipeline consists of:

\begin{enumerate}
    \item \textbf{Candidate Generation:} Producing a set of candidate items (movies) for each user.
    \item \textbf{Ranker:} Scoring candidates based on predicted relevance to the user.
    \item \textbf{Reranker:} Applying list-level constraints (diversity) to the ranked candidates.
    \item \textbf{Personalization Update:} Updating the user's state based on simulated feedback.
\end{enumerate}

This report documents the implementation details of each component, presents empirical results from required experiments, and provides answers to theoretical questions in the appendix. All experiments are designed to be reproducible, with hyperparameters logged and a fast mode implemented for efficient prototyping.

\section{Dataset and Preprocessing}
\label{sec:dataset}

We use the MovieLens 100K dataset, which contains 100,000 ratings from 943 users on 1,682 movies. Each rating is on a scale of 1 to 5.

\subsection{Train/Test Split}
We performed an 80/20 random split over observed ratings. Users with no test items rated $\geq 4$ were excluded from evaluation for ranking metrics. For candidate generation, we excluded all movies rated by the user in the training set to avoid recommending previously interacted items.

\subsection{Relevance Definition}
For ranking evaluation, we defined relevant items for user $u$ as:
\[
G_u = \{i : (u,i) \in \Omega_{\text{test}} \text{ and } r_{ui} \geq 4\}
\]

\subsection{Feature Engineering}
For the Pairwise LTR ranker, we constructed feature vectors $f(u,i)$ containing:
\begin{itemize}
    \item $s_{\text{MF}}(u,i)$: Predicted score from a pre-trained MF model
    \item $b_u$: User bias term
    \item $b_i$: Item bias term
    \item $\text{pop}(i)$: Global popularity of item $i$ (log-transformed interaction count)
    \item 1: Intercept term
\end{itemize}

For diversity calculations and MMR similarity, we used genre vectors derived from the `u.item` file, representing each movie as a one-hot vector over 19 genres.

\section{Implementation Details}
\label{sec:implementation}

\subsection{Rankers}
We implemented four ranking methods:

\subsubsection*{R1: Popularity Baseline}
Movies are ranked globally by their number of interactions in the training set. This serves as a non-personalized baseline.

\subsubsection*{R2: MF-SGD}
Matrix Factorization with biases was implemented with the following prediction function:
\[
\hat{r}_{ui} = \mu + b_u + b_i + p_u^\top q_i
\]
Parameters were optimized using Stochastic Gradient Descent with squared error loss and L2 regularization. Hyperparameters: $d=20$, $\lambda=0.05$, $\eta=0.01$, epochs $=10$.

\subsubsection*{R3: MF-ALS}
The same MF model was optimized using Alternating Least Squares, which solves for user and item factors in closed form with ridge regularization. We used $\lambda=0.05$ for 8 iterations.

\subsubsection*{R4: Pairwise LTR}
We trained a Bradley-Terry logistic pairwise model with features $f(u,i)$:
\[
s_{\phi}(u,i) = \phi^\top f(u,i)
\]
Training used positive $(i^+)$ and negative $(i^-)$ item pairs per user, optimizing:
\[
\mathcal{L} = -\sum_{(u,i^+,i^-)} \log \sigma(s_{\phi}(u,i^+) - s_{\phi}(u,i^-))
\]
We sampled up to 50 pairs per user and trained for 6 epochs.

\subsection{Reranking: MMR}
For the MF baseline, we applied Maximal Marginal Relevance reranking over the top $M=80$ candidates to produce a final slate of size $K=10$:
\[
c^{\star} = \underset{c \in C \setminus S}{\arg\max}\left((1 - \alpha)\ \text{rel}(c) - \alpha \max_{s \in S} \text{sim}(c,s)\right)
\]
We used $\alpha \in \{0.1, 0.4, 0.7\}$ and genre-based cosine similarity.

\subsection{Personalization Update}
To simulate session-based personalization, we updated the user's latent factor using Exponential Moving Average:
\[
u_{t+1} = \frac{(1 - \rho) u_t + \rho v_{i_t}}{\|(1 - \rho) u_t + \rho v_{i_t}\|}
\]
where $v_{i}$ is the learned item factor from MF, and $\rho \in \{0.1, 0.3\}$. For demonstration, we simulated a user selecting the top-ranked item from their initial slate.

\section{Experiments and Results}
\label{sec:experiments}

\subsection{E1: Baseline Comparison}
\label{sec:e1}

We compared the four ranking methods using Recall@10, NDCG@10, and Diversity@10 averaged over users with at least one relevant test item. Results are reported in Table \ref{tab:baseline_comparison}.

\begin{table}[H]
\centering
\caption{Comparison of Ranking Methods}
\label{tab:baseline_comparison}
\begin{tabular}{lccc}
\toprule
\textbf{Method} & \textbf{Recall@10} & \textbf{NDCG@10} & \textbf{Diversity@10} \\
\midrule
Popularity Baseline & [Value] & [Value] & [Value] \\
MF-SGD & [Value] & [Value] & [Value] \\
MF-ALS & [Value] & [Value] & [Value] \\
Pairwise LTR & [Value] & [Value] & [Value] \\
\bottomrule
\end{tabular}
\end{table}

\textbf{Analysis:} 
[Discuss the results. Which method performed best? How do MF-SGD and MF-ALS compare? Does the pairwise LTR model improve ranking quality?]

\subsection{E2: MMR Tradeoff}
\label{sec:e2}

We evaluated the effect of MMR reranking on the MF-SGD baseline across different $\alpha$ values. Figure \ref{fig:mmr_tradeoff} shows NDCG@10 versus Diversity@10.

\begin{figure}[H]
\centering
\includegraphics[width=0.7\textwidth]{placeholder_plot.png}
\caption{NDCG@10 vs Diversity@10 for MMR with varying $\alpha$}
\label{fig:mmr_tradeoff}
\end{figure}

\textbf{Analysis:}
[Describe the tradeoff observed. As $\alpha$ increases, what happens to relevance (NDCG) and diversity? Which $\alpha$ provides a good balance?]

\subsection{E3: Ablation Study}
\label{sec:e3}

We conducted ablation experiments to understand the contribution of different components:

\subsubsection*{Effect of MMR}
We compared MF-SGD with and without MMR ($\alpha=0.4$) as shown in Table \ref{tab:mmr_ablation}.

\begin{table}[H]
\centering
\caption{Ablation: Effect of MMR}
\label{tab:mmr_ablation}
\begin{tabular}{lccc}
\toprule
\textbf{Method} & \textbf{Recall@10} & \textbf{NDCG@10} & \textbf{Diversity@10} \\
\midrule
MF-SGD (no MMR) & [Value] & [Value] & [Value] \\
MF-SGD + MMR ($\alpha=0.4$) & [Value] & [Value] & [Value] \\
\bottomrule
\end{tabular}
\end{table}

\subsubsection*{Effect of Popularity Feature}
We trained the Pairwise LTR model with and without the popularity feature to assess its impact. Results are in Table \ref{tab:pop_ablation}.

\begin{table}[H]
\centering
\caption{Ablation: Effect of Popularity Feature}
\label{tab:pop_ablation}
\begin{tabular}{lccc}
\toprule
\textbf{Method} & \textbf{Recall@10} & \textbf{NDCG@10} & \textbf{Diversity@10} \\
\midrule
Pairwise LTR (full features) & [Value] & [Value] & [Value] \\
Pairwise LTR (no pop) & [Value] & [Value] & [Value] \\
\bottomrule
\end{tabular}
\end{table}

\textbf{Analysis:}
[Discuss the impact of removing MMR. How does it affect relevance vs diversity? Discuss the importance of the popularity feature in the LTR model. Does it help or hurt performance?]

\subsection{Personalization Demonstration}
We simulated a user session using EMA updates. Table \ref{tab:ema_update} shows the change in user state and subsequent recommendations after selecting a movie.

\begin{table}[H]
\centering
\caption{EMA Personalization Update ($\rho=0.3$)}
\label{tab:ema_update}
\begin{tabular}{ll}
\toprule
\textbf{Stage} & \textbf{Top-3 Recommended Items} \\
\midrule
Initial & [Item A], [Item B], [Item C] \\
After selecting [Item X] & [Item D], [Item E], [Item F] \\
\bottomrule
\end{tabular}
\end{table}

\textbf{Analysis:}
[Describe how the user state changed and how recommendations shifted. Compare results for $\rho=0.1$ vs $\rho=0.3$.]

\section{Reproducibility}
\label{sec:reproducibility}

\subsection{Fast Mode}
Our implementation includes a `FAST_MODE` flag that:
\begin{itemize}
    \item Sets $d \leq 20$
    \item Limits SGD epochs to $\leq 10$ and ALS iterations to $\leq 8$
    \item Evaluates on a subset of 300 eligible users
    \item Completes end-to-end in under 10 minutes on Colab free tier
\end{itemize}

\subsection{JSONL Logging}
We produce a JSONL log file where each line corresponds to one evaluation request, containing:
\begin{itemize}
    \item user id (internal index)
    \item method name
    \item top-10 list of item indices
    \item metrics (NDCG@10, Recall@10, Diversity@10)
    \item hyperparameters ($d$, $\lambda$, $\eta$, $\alpha$, $\rho$)
    \item (for sessions) chosen item and updated state summary
\end{itemize}

\section{Conclusion}
\label{sec:conclusion}

[Summarize the main findings of your project. What did you learn about the tradeoffs between relevance and diversity? Which ranking method performed best? Reflect on the challenges of offline evaluation and the importance of co-creative systems.]

\begin{appendices}
\section{Theory Questions}
\label{app:theory}

\subsection*{T1: ALS Closed-Form Update}
Given fixed $Q$, biases, and residual $e_{ui} = r_{ui} - (\mu + b_u + b_i)$, we optimize:
\[
\min_{p_u} \sum_{i \in \Omega(u)} (e_{ui} - p_u^\top q_i)^2 + \lambda \|p_u\|^2
\]
This is a ridge regression problem. Taking the gradient with respect to $p_u$ and setting to zero:
\[
\nabla_{p_u} = -2\sum_{i \in \Omega(u)} (e_{ui} - p_u^\top q_i) q_i + 2\lambda p_u = 0
\]
\[
\sum_{i \in \Omega(u)} (e_{ui} - p_u^\top q_i) q_i = \lambda p_u
\]
\[
\sum_{i \in \Omega(u)} e_{ui} q_i - \sum_{i \in \Omega(u)} q_i q_i^\top p_u = \lambda p_u
\]
\[
\sum_{i \in \Omega(u)} e_{ui} q_i = \left( \sum_{i \in \Omega(u)} q_i q_i^\top + \lambda I \right) p_u
\]
Thus:
\[
p_u = \left( \sum_{i \in \Omega(u)} q_i q_i^\top + \lambda I \right)^{-1} \left( \sum_{i \in \Omega(u)} e_{ui} q_i \right)
\]

\subsection*{T2: SGD Updates for MF with Biases}
For prediction $\hat{r}_{ui} = \mu + b_u + b_i + p_u^\top q_i$ with loss $\mathcal{L} = \frac{1}{2}(r_{ui} - \hat{r}_{ui})^2 + \frac{\lambda}{2}(\|p_u\|^2 + \|q_i\|^2 + b_u^2 + b_i^2)$, the gradients are:
\[
\frac{\partial \mathcal{L}}{\partial p_u} = -e_{ui} q_i + \lambda p_u
\]
\[
\frac{\partial \mathcal{L}}{\partial q_i} = -e_{ui} p_u + \lambda q_i
\]
\[
\frac{\partial \mathcal{L}}{\partial b_u} = -e_{ui} + \lambda b_u
\]
\[
\frac{\partial \mathcal{L}}{\partial b_i} = -e_{ui} + \lambda b_i
\]
where $e_{ui} = r_{ui} - \hat{r}_{ui}$. The SGD updates are:
\[
p_u \leftarrow p_u + \eta (e_{ui} q_i - \lambda p_u)
\]
\[
q_i \leftarrow q_i + \eta (e_{ui} p_u - \lambda q_i)
\]
\[
b_u \leftarrow b_u + \eta (e_{ui} - \lambda b_u)
\]
\[
b_i \leftarrow b_i + \eta (e_{ui} - \lambda b_i)
\]

\subsection*{T3: Pairwise Gradient}
Given $s_{\phi}(u,i) = \phi^\top f(u,i)$ and $\Delta f = f(u,i^+) - f(u,i^-)$, we have:
\[
\mathcal{L} = -\log \sigma(\phi^\top \Delta f) = -\log\left(\frac{1}{1 + e^{-\phi^\top \Delta f}}\right)
\]
\[
\frac{\partial \mathcal{L}}{\partial \phi} = -\frac{\partial}{\partial \phi} \log \sigma(\phi^\top \Delta f)
\]
Using $\sigma'(x) = \sigma(x)(1 - \sigma(x))$:
\[
\nabla_{\phi} \mathcal{L} = -\frac{1}{\sigma(\phi^\top \Delta f)} \cdot \sigma(\phi^\top \Delta f)(1 - \sigma(\phi^\top \Delta f)) \cdot \Delta f
\]
\[
\nabla_{\phi} \mathcal{L} = -(1 - \sigma(\phi^\top \Delta f)) \Delta f
\]
Thus:
\[
\nabla_{\phi} \left( -\log \sigma(\phi^\top \Delta f) \right) = (\sigma(\phi^\top \Delta f) - 1) \Delta f
\]

\subsection*{T4: Why RMSE $\neq$ Ranking Quality?}
Optimizing RMSE focuses on accurately predicting the exact rating value for all items, but ranking quality depends only on the relative order of items. A model can have low RMSE but poor top-K ranking if:
\begin{itemize}
    \item Prediction errors are not uniform across items; systematic biases can shuffle rankings.
    \item RMSE weights all items equally, but ranking metrics (e.g., NDCG) focus on top positions.
    \item User preferences are ordinal; a model that predicts 3.5 vs 3.6 correctly may still misorder items if the true gap is small.
\end{itemize}
Thus, ranking requires optimizing for relative ordering, not absolute value accuracy.

\subsection*{T5: Offline Evaluation Caveat}
Offline evaluation assumes that the data reflects unbiased user preferences. However, when exposure is not random:
\begin{itemize}
    \item We only observe ratings for items that were shown to users (exposure bias).
    \item Popular items are overrepresented in training data.
    \item A model that simply recommends popular items may appear good offline but fails to discover niche items users might prefer.
    \item Without randomization, we cannot distinguish whether a high rating is due to user preference or because the item was prominently displayed.
\end{itemize}
Thus, offline metrics may overestimate real-world performance, especially for novel or long-tail items.

\end{appendices}

\end{document}